% Introduction to Part V. Unnumbered text: no \label except the \phantomsection target.
\phantomsection\label{intro:V}
Part~V solves Ba\l aban's variational problem relative to a given reference: minimize the Wilson action on the set $\mathfrak B(V)$ of fine configurations of small curvature whose $k$-fold average is a given real coarse configuration $V$ (Section~\ref{sec:var-theorem}). The reference $U_0$ is real, $(M,\eps)$-regular, with current $|J_{U_0}|_\infty\le\eps$, and the averages of $V$ are close to those of $U_0$: $V(c)W_k(c)^{-1}=e^{\ii B(c)}$ with $|B|_\infty\le\beta$. The critical points are Ba\l aban's background fields. Part~VI supplies the reference by an induction over $k$, and Part~VII expands the renormalization-group step around the background field, with its multiplier as the gradient of the coarse action.

\paragraph{What this Part proves.} The main result is Theorem~\ref{thm:variational}, which is Theorem~\ref{thm:I}. If $\eps+\beta$ is small, then:
\begin{itemize}
\item the action on $\mathfrak B(V)$ has exactly one critical $\mathfrak G_0$-orbit among configurations of small curvature and current, and it consists of the minimizers there (Theorem~\ref{thm:variational}(a),(b));
\item its Landau representative $A_*$ relative to $U_0$ is small with its derivatives, and $A_*=\bH B-\bC J_{U_0}+O((\eps+\beta)^2)$ (Theorem~\ref{thm:variational}(c));
\item there is a unique Euler--Lagrange multiplier $\omega$ on the unit lattice (Theorem~\ref{thm:variational}(d));
\item the minimizer is $(M,\eps+C_6(\eps+\beta))$-regular, $A_*$ is analytic in complex $B$, and $dA_*/dB$ decays exponentially (Theorem~\ref{thm:variational}(e),(f));
\item the orbit is covariant under gauge transformations and under the lattice symmetries $\mathfrak S_k$ (Theorem~\ref{thm:variational}(g),(h)).
\end{itemize}
The constants do not depend on $k$, $n$ or $\eta$; they depend on $L$, and may be taken independent of $M\ge M_0$. The proof has four steps (Section~\ref{sec:var-theorem}): reduction to the Landau chart of Part~IV; a linearizing transformation of the constraint and a contraction for $A_1=-\bC(J+\nabla\mathcal V(A_1+\bH B))$ (Lemma~\ref{lem:linearizing}, Proposition~\ref{prop:contraction}); convexity, which gives minimality (Proposition~\ref{prop:minimality}); and the Euler--Lagrange equation, through the gauge covariance of the gauge-adapted averaging (Lemma~\ref{lem:Q-covariance}, Proposition~\ref{prop:euler-lagrange}).

\paragraph{Imports.} From Part~IV: Theorem~\ref{thm:chart}, Corollary~\ref{cor:parametrization}, Propositions~\ref{prop:Q-expansion}, \ref{prop:linearized-averages} and~\ref{prop:action}, and Lemma~\ref{lem:sup-bounds}. Parts~II and~III enter only through the pointwise bounds of Lemma~\ref{lem:sup-bounds}, the identities for $\bH$ and $\bC$ (Theorem~\ref{thm:cov-min}) and the coercivity on the tangent space (Lemma~\ref{lem:X-coercive}); the random-walk expansion is not used (Remark~\ref{x:27.6}). No proof in Part~V uses a later result. The constants are fixed at the end of Section~\ref{sec:var-regularity}; one condition there serves only the Dirichlet version of the contraction in Section~\ref{sec:interaction-local}, which is off the main line.

\paragraph{Exports.} Theorem~\ref{thm:variational} goes to the induction (Lemma~\ref{lem:improvement}, Theorem~\ref{thm:balaban}), to Part~VII, and to Parts~IX--XI. The linearizing transformation, the bounds on $\mathcal V$ and the contraction (Lemmas~\ref{lem:linearizing}, \ref{lem:V-bounds}, Proposition~\ref{prop:contraction}) are used again for the coarse action in Part~VII and for the weakened background map in Part~IX (Section~\ref{sec:weak-background}). The Euler--Lagrange equation (Proposition~\ref{prop:euler-lagrange}) goes to Parts~VI, VII and~XI. Lemma~\ref{lem:Lstar-coercive}, Proposition~\ref{prop:minimality} and Lemma~\ref{lem:reduction} are used again only in Part~XI.

\paragraph{Details for this Part.} The companion has no chapter on Part~V; its parameter ledger lists the constants of Part~V that Parts~VII--XI use, with their order of choice.
